\section{Exact sharpness of the geometric projection step}
\label{sec:geometricsharp}
The preceding family proves sharpness for the original invariant.
Here we separately verify sharpness of the existence-of-direction bound
in Lemma~\ref{lem:collapse}. Equality in just one direction would not
suffice for this purpose; an upper bound in every direction is needed.

Fix integers $1\le r\le d-1$, put $m=r+1$, and let $0<\rho<1$. Set
\begin{align*}
 P&=\conv(0,e_1,\ldots,e_d),\\
 K&=\conv(e_1,\ldots,e_d,\rho e_1,\ldots,\rho e_m)\\
  &=\left\{x\ge0:\sum_i x_i\le1,\quad
       \rho^{-1}\sum_{i\le m}x_i+\sum_{i>m}x_i\ge1\right\}.
\end{align*}
The barycentric loss at zero is $\rho$. This body has $d+m=d+1+r$
vertices and is an iterated pyramid over the $m$-dimensional equal
truncation. Every facet plane of $P$ supports a full facet of $K$.

With $c=1/(d-1)!$, the actual facet area-normal vectors of $K$ are
\begin{align*}
 &-c(1-\rho^{m-1})e_i &&(i\le m),\\
 &-c(1-\rho^m)e_i &&(i>m),\\
 &c\mathbf1,\qquad
 -c(\underbrace{\rho^{m-1},\ldots,\rho^{m-1}}_{m},
       \underbrace{\rho^m,\ldots,\rho^m}_{d-m}).
\end{align*}
For a coordinate facet the removed simplex has the corresponding product
of intercept scales; the last vector is obtained from the intercepts of
the new cut facet. These vectors sum to zero, as required.

For nonzero $z\in\R^d$ write
\[
 A=\sum_{i\le m}z_i,\quad B=\sum_{i>m}z_i,\quad
 A_* =\sum_{i\le m}|z_i|,\quad B_* =\sum_{i>m}|z_i|.
\]
Cauchy's formula gives the exact \emph{homogeneous} brightness ratio
\begin{equation}\label{eq:sharpbrightness}
 \frac{b_P(z)-b_K(z)}{b_P(z)}
 =\rho^{m-1}\frac{A_*+\rho B_*-|A+\rho B|}
                         {A_*+B_*+|A+B|}.
\end{equation}
The denominator is positive. The numerator is nonnegative by the triangle
inequality and is at most the denominator. Therefore the relative deficit
is at most $\rho^{m-1}=\rho^r$ in every direction. For
$z=e_1-e_2$, the ratio equals $\rho^r$, and that direction is parallel to
the opposite base facet of $P$. Upon normalizing $z$, we conclude
\begin{equation}\label{eq:maxsharp}
 \max_{\norm u=1}\frac{\pi_P(u)-\pi_K(u)}{\pi_P(u)}=\rho^r.
\end{equation}
The same maximum is attained among base-parallel unit directions.
Thus both the exponent and the constant one in
Lemma~\ref{lem:collapse} are sharp throughout $1\le r\le d-1$.
The distinction between $b_K(z)$ and $\pi_K(u)$ avoids any convention for
projection notation at nonunit vectors.
